% Archivo compartido: no compilar directamente.
% Use clase_1_estudiante.tex o clase_1_docente.tex como archivo raíz.
\ifdefined\documentoSemanal
\else
\usepackage{tikz}
\usepackage{float}
\usepackage{estilo_apuntes}
\usetikzlibrary{arrows.meta,positioning,babel}

\encabezado{Semana 1 · Clase 1}{Material de clase}
\begin{document}
\fi
\tituloclase{Funciones: conceptos fundamentales}{Función, dominio, imagen y restricción}{Semana 1 · Clase 1 · 1 hora}

\begin{objetivos}
\begin{itemize}
  \item Reconocer cuándo una relación define una función.
  \item Distinguir dominio, codominio e imagen.
  \item Interpretar la notación $f:A\to B$ y el valor $f(x)$.
  \item Comprender la restricción de una función.
\end{itemize}
\end{objetivos}      

\begin{resumen}
Una función es una relación especial entre dos conjuntos: cada elemento del conjunto de partida debe estar relacionado con exactamente un elemento del conjunto de llegada. La función queda determinada por su dominio, su codominio y su regla de asignación.
\end{resumen}

\section*{Relaciones y funciones}

Dados dos conjuntos $A$ y $B$, el producto cartesiano es
\[
A\times B=\{(x,y):x\in A,\ y\in B\}.
\]
Una relación de $A$ en $B$ es cualquier subconjunto de $A\times B$.

\begin{definicion}[función]
Una relación $f\subseteq A\times B$ es una función de $A$ en $B$ si:
\begin{enumerate}
  \item para todo $x\in A$ existe $y\in B$ tal que $(x,y)\in f$;
  \item si $(x,y)\in f$ y $(x,z)\in f$, entonces $y=z$.
\end{enumerate}
Se escribe $f:A\to B$ y $(x,y)\in f$ se expresa como $f(x)=y$.
\end{definicion}

\begin{nota}
En otras palabras, $f$ es una función de $A$ en $B$ si es una relación entre los
elementos de $A$ y $B$ de modo que para cada elemento $x$ de $A$, hay un único
elemento $y$ de $B$ que le corresponde a $x$ en esta relación; a ese elemento y
se le llama imagen de $x$ respecto de $f$ y se le representa por $f(x)$.
\end{nota}

La primera condición garantiza que ninguna entrada queda sin imagen. La segunda garantiza que la imagen de cada entrada sea única. Entradas diferentes pueden compartir una misma salida.

\begin{ejemplo}[función y contraejemplo]
Determine si las siguientes relaciones 
\[
R=\{(1,2),(2,2),(3,5)\},\qquad
S=\{(1,2),(1,4),(3,5)\}
\]
definen funciones desde sus conjuntos de primeras componentes hacia $\mathbb{R}$.
\end{ejemplo}
\ifmostrarSoluciones
\begin{solucion}
$R$ sí es función: cada entrada $1,2,3$ aparece asociada con una única salida. Que $1$ y $2$ compartan la salida $2$ no viola la definición. $S$ no es función porque la entrada $1$ posee dos salidas distintas.
\end{solucion}
\else\vspace{12mm}\fi

\begin{figure}[H]
\centering
\begin{tikzpicture}[x=1.05cm,y=0.62cm,>=Stealth,every node/.style={font=\small}]
  \draw[->,Gris] (-0.35,0) -- (4.15,0) node[right] {$x$};
  \draw[->,Gris] (0,-0.45) -- (0,6.15) node[above] {$y$};
  \foreach \x in {1,2,3}{
    \draw[Gris] (\x,0.10)--(\x,-0.10) node[below] {\x};
  }
  \foreach \y in {1,2,3,4,5}{
    \draw[Gris] (0.10,\y)--(-0.10,\y) node[left] {\y};
  }
  \draw[dashed,Azul!55] (1,0)--(1,2)--(0,2);
  \draw[dashed,Azul!55] (2,0)--(2,2);
  \draw[dashed,Azul!55] (3,0)--(3,5)--(0,5);
  \fill[Azul] (1,2) circle (2.4pt) node[above left] {$(1,2)$};
  \fill[Azul] (2,2) circle (2.4pt) node[above right] {$(2,2)$};
  \fill[Azul] (3,5) circle (2.4pt) node[above right] {$(3,5)$};
\end{tikzpicture}
\caption{Gráfico de $R=\{(1,2),(2,2),(3,5)\}$. Cada valor de $x$ aparece con una única coordenada $y$, por lo que $R$ representa una función.}
\label{fig:pares-ordenados-funcion}
\end{figure}

\begin{definicion}[dominio y codominio]
En $f:A\to B$, el conjunto $A$ es el \textbf{dominio} y $B$ es el \textbf{codominio}.
\end{definicion}

\begin{definicion}[imagen o recorrido]
La imagen de $f:A\to B$ es
\[
\operatorname{Im}(f)=\{f(x):x\in A\}\subseteq B.
\]
Contiene únicamente las salidas alcanzadas y no tiene que coincidir con el codominio.
\end{definicion}

\begin{figure}[H]
\centering
\begin{tikzpicture}[>=Stealth,thick,every node/.style={font=\small}]
  \draw[Azul,fill=AzulClaro] (0,0) ellipse (1.25cm and 2.05cm);
  \draw[Verde,fill=VerdeClaro] (5.2,0) ellipse (1.25cm and 2.05cm);
  \node[Azul,font=\bfseries] at (0,2.35) {$A$ (dominio)};
  \node[Verde,font=\bfseries] at (5.2,2.35) {$B$ (codominio)};
  \node (x1) at (0,1.15) {$x_1$};
  \node (x2) at (0,0) {$x_2$};
  \node (x3) at (0,-1.15) {$x_3$};
  \node (y1) at (5.2,1.15) {$y_1$};
  \node (y2) at (5.2,0) {$y_2$};
  \node (y3) at (5.2,-1.15) {$y_3$};
  \draw[->,Azul] (x1) -- (y1);
  \draw[->,Azul] (x2) -- (y1);
  \draw[->,Azul] (x3) -- (y3);
  \node[align=center,text=Gris,font=\footnotesize] at (2.6,-2.0)
    {cada elemento de $A$ tiene\\exactamente una flecha de salida};
\end{tikzpicture}
\caption{Representación de una función $f:A\to B$. Dos entradas distintas pueden compartir una imagen; $y_2$ puede no pertenecer a $\operatorname{Im}(f)$.}
\label{fig:diagrama-funcion}
\end{figure}

\begin{ejemplo}[dominio, codominio e imagen]
Sea $f:\{-2,-1,0,1,2\}\to\mathbb{R}$, $f(x)=x^2$. Identifique su dominio, su codominio y su imagen.
\end{ejemplo}
\ifmostrarSoluciones
\begin{solucion}
\[
\operatorname{Dom}(f)=\{-2,-1,0,1,2\},\qquad B=\mathbb{R}.
\]
Las salidas obtenidas son $4,1,0,1,4$, por lo que $\operatorname{Im}(f)=\{0,1,4\}$.
\end{solucion}
\else\vspace{12mm}\fi

\begin{definicion}[gráfica de $f$]
El gráfico de una relación $f:A\to\mathbb{R}$ es
\[
G_f=\{(x,f(x)):x\in A\}.
\]
\end{definicion}
\begin{proposicion}[de la línea vertical]
\begin{itemize}
\item Dada una función todas las líneas verticales que se puedan dibujar intersecan el gráfico de $f$ solo una vez. 
\item Si cualquier línea vertical interseca al gráfico de una relación $f$ más de una vez, entonces esta relación $f$ no representa una función.
\end{itemize}
\end{proposicion}

\begin{figure}[H]
\centering
\begin{tikzpicture}[>=Stealth,every node/.style={font=\footnotesize}]
  % Panel A: gráfico de una función.
  \begin{scope}[xshift=0cm,yshift=-1.25cm,x=0.72cm,y=0.62cm]
    \draw[->,Gris] (-2.45,0)--(2.45,0) node[right] {$x$};
    \draw[->,Gris] (0,-0.45)--(0,4.6) node[above] {$y$};
    \draw[Azul,very thick,domain=-2.05:2.05,samples=80,smooth,variable=\x]
      plot ({\x},{\x*\x});
    \draw[Rojo,dashed,thick] (1,-0.25)--(1,4.35)
      node[above,text=Rojo] {$x=1$};
    \fill[Rojo] (1,1) circle (2.3pt);
    \node[Azul,font=\bfseries] at (0,-1.0) {(a) Sí es función};
    \node[align=center,text=Gris] at (0,-1.55)
      {una sola intersección};
  \end{scope}

  % Panel B: relación que no es función de x.
  \begin{scope}[xshift=5.6cm,x=0.72cm,y=0.62cm]
    \draw[->,Gris] (-2.45,0)--(2.45,0) node[right] {$x$};
    \draw[->,Gris] (0,-2.45)--(0,2.45) node[above] {$y$};
    \draw[Azul,very thick] (0,0) circle (2);
    \draw[Rojo,dashed,thick] (1,-2.35)--(1,2.35)
      node[above,text=Rojo] {$x=1$};
    \fill[Rojo] (1,{sqrt(3)}) circle (2.3pt);
    \fill[Rojo] (1,{-sqrt(3)}) circle (2.3pt);
    \node[Azul,font=\bfseries] at (0,-3.0) {(b) No es función};
    \node[align=center,text=Gris] at (0,-3.55)
      {dos intersecciones};
  \end{scope}
\end{tikzpicture}
\caption{Criterio de la línea vertical. En (a), cualquier recta vertical corta el gráfico a lo sumo una vez. En (b), al menos una recta vertical corta la circunferencia dos veces.}
\label{fig:criterio-linea-vertical}
\end{figure}


\begin{definicion}[restricción]
Si $C\subseteq A$, la restricción de $f:A\to B$ a $C$ es
\[
f|_C:C\to B,\qquad f|_C(x)=f(x).
\]
La regla se conserva, pero se reduce el dominio y, posiblemente, la imagen.
\end{definicion}


\begin{ejemplo}[efecto de una restricción]
Compare $q:\mathbb{R}\to\mathbb{R}$, $q(x)=x^2$, con $q|_{[0,\infty)}$.
\end{ejemplo}
\ifmostrarSoluciones
\begin{solucion}
La regla es la misma. Los dominios son $\mathbb{R}$ y $[0,\infty)$, respectivamente; ambas imágenes son $[0,\infty)$. En la restricción desaparecen los pares de entradas opuestas con una misma salida.
\end{solucion}
\else\vspace{12mm}\fi

\begin{figure}[H]
\centering
\begin{tikzpicture}[>=Stealth,every node/.style={font=\footnotesize}]
  % Función sobre todo R.
  \begin{scope}[xshift=0cm,x=0.72cm,y=0.58cm]
    \draw[->,Gris] (-2.65,0)--(2.65,0) node[right] {$x$};
    \draw[->,Gris] (0,-0.35)--(0,5.25) node[above] {$y$};
    \foreach \x in {-2,-1,1,2}{
      \draw[Gris] (\x,0.09)--(\x,-0.09) node[below] {\x};
    }
    \foreach \y in {1,4}{
      \draw[Gris] (0.09,\y)--(-0.09,\y) node[left] {\y};
    }
    \draw[Azul,very thick,domain=-2.25:2.25,samples=80,smooth,variable=\x]
      plot ({\x},{\x*\x});
    \node[Azul,font=\bfseries] at (0,-1.0) {$q:\mathbb{R}\to\mathbb{R}$};
    \node[Azul] at (0,-1.55) {$q(x)=x^2$};
  \end{scope}

  % Restricción a [0,infinito).
  \begin{scope}[xshift=5.6cm,x=0.72cm,y=0.58cm]
    \draw[->,Gris] (-0.55,0)--(2.65,0) node[right] {$x$};
    \draw[->,Gris] (0,-0.35)--(0,5.25) node[above] {$y$};
    \foreach \x in {1,2}{
      \draw[Gris] (\x,0.09)--(\x,-0.09) node[below] {\x};
    }
    \foreach \y in {1,4}{
      \draw[Gris] (0.09,\y)--(-0.09,\y) node[left] {\y};
    }
    \draw[Verde,very thick,domain=0:2.25,samples=60,smooth,variable=\x]
      plot ({\x},{\x*\x});
    \fill[Verde] (0,0) circle (2.2pt);
    \node[Verde,font=\bfseries] at (1.0,-1.0) {$q|_{[0,\infty)}$};
    \node[Verde] at (1.0,-1.55) {$q|_{[0,\infty)}(x)=x^2$};
  \end{scope}
\end{tikzpicture}
\caption{La restricción conserva la regla $x^2$ y únicamente la rama correspondiente a $x\geq0$. Ambas funciones tienen imagen $[0,\infty)$, pero dominios diferentes.}
\label{fig:restriccion-cuadratica}
\end{figure}

\begin{ejercicios}[para la clase]
\begin{enumerate}
  \item Determine cuáles relaciones definen una función:
  \[
  A=\{(-1,0),(0,1),(1,0)\},\quad
  B=\{(-1,0),(-1,1),(1,2)\}.
  \]
  \item Sea $f:\mathbb{R}\to\mathbb{R}$, $f(x)=2x^2-x+3$. Calcule $f(0)$, $f(-2)$, $f(a)$ y $f(a+h)$.
  \item Si $g:\{-2,-1,0,1,2\}\to\mathbb{Z}$, $g(x)=|x|-1$, encuentre su imagen.
  \item Distinga dominio, codominio e imagen para $h:\mathbb{R}\to\mathbb{R}$, $h(x)=x^2+1$.
  %\item Proponga una función con codominio $\mathbb{R}$ e imagen $[-1,1]$.
  \item Explique por qué $x^2+y^2=4$ no define a $y$ como función de $x$ y obtenga dos restricciones que si sean funciones.
  %\item Sea $p:\mathbb{R}\to\mathbb{R}$, $p(x)=x^2-4$. Identifique el dominio y la imagen de $p|_{[-1,3]}$.
  \item Construya una función definida en $\{1,2,3,4\}$ cuya imagen tenga exactamente dos elementos.
\end{enumerate}
\end{ejercicios}

\ifmostrarSoluciones
\begin{solucion}[de los ejercicios]
\begin{enumerate}
  \item $A$ es función; $B$ no lo es porque $-1$ posee dos imágenes.
  \item $f(0)=3$, $f(-2)=13$, $f(a)=2a^2-a+3$ y
  \[
  f(a+h)=2a^2+4ah+2h^2-a-h+3.
  \]
  \item $\operatorname{Im}(g)=\{-1,0,1\}$.
  \item Dominio $\mathbb{R}$, codominio $\mathbb{R}$ e imagen $[1,\infty)$.
  %\item Una posibilidad es $f(x)=\sin x$ con dominio y codominio $\mathbb{R}$.
  \item $y=\pm\sqrt{4-x^2}$, $-2\leq x\leq2$, define los semicírculos superior e inferior.
  %\item Dominio $[-1,3]$ e imagen $[-4,5]$.
  \item Por ejemplo, $f(1)=f(2)=0$ y $f(3)=f(4)=1$.
\end{enumerate}
\end{solucion}
\fi

\fuentes{\textbf{Para ampliar:} J. Stewart, \emph{Cálculo de una variable: trascendentes tempranas}, §1.1, y G. Strang--E. Herman, \emph{Cálculo, volumen 1}, OpenStax, §1.1.}
\ifdefined\documentoSemanal
\else
\end{document}
\fi
